% Part: first-order-logic
% Chapter: sequent-calculus

\documentclass[../../../include/open-logic-chapter]{subfiles}

\begin{document}

\iftag{FOL}
      {\olchapter{fol}{seq}{সিকোয়েন্ট কলন}}
      {\olchapter{pl}{seq}{সিকোয়েন্ট কলন}}

\begin{editorial}
  এই অধ্যায়ে ধ্রুপদি প্রথম-ক্রমের যুক্তিবিদ্যার জন্য গেন্‌ৎসেনের প্রমিত
  সিকোয়েন্ট কলন LK উপস্থাপন করা হয়েছে। আরও উদাহরণ ও অনুশীলনী
  যোগ করা যেতে পারে। প্রমাণপদ্ধতি হিসেবে সিকোয়েন্ট কলনের প্রাসঙ্গিক
  উপকরণ অন্তর্ভুক্ত বা বাদ দিতে ``prfLK'' ট্যাগটি ব্যবহার করুন।
\end{editorial}

\olimport{rules-and-proofs}

\olimport{propositional-rules}

\iftag{FOL}{%
  \olimport{quantifier-rules}
}{}

\olimport{structural-rules}

\olimport{derivations}

\olimport{proving-things}

\iftag{FOL}{%
  \olimport{proving-things-quant}
}{}

\olimport{proof-theoretic-notions}

\olimport{provability-consistency}

\olimport{provability-propositional}

\iftag{FOL}{%
  \olimport{provability-quantifiers}
}{}

\olimport{soundness}

\iftag{FOL}{%
  \olimport{identity}
  \olimport{soundness-identity}
}{}

\OLEndChapterHook

\end{document}
